\documentclass[11pt]{article}
% jugeo-common.tex — shared preamble for all Judgment Geometry papers
% Include via: % jugeo-common.tex — shared preamble for all Judgment Geometry papers
% Include via: % jugeo-common.tex — shared preamble for all Judgment Geometry papers
% Include via: \input{jugeo-common}

% ── Packages ────────────────────────────────────────────────────────────
\usepackage{amsmath,amssymb,amsthm,mathtools}
\usepackage{tikz-cd}
\usepackage{listings}
\usepackage{xcolor}
\usepackage{xspace}
\usepackage{booktabs}
\usepackage{multirow}
\usepackage{enumitem}
\usepackage[expansion=false]{microtype}
\usepackage{hyperref}
\usepackage{cleveref}
% \usepackage{thmtools}  % disabled: not available in this TeX installation
\usepackage{float}
\usepackage{caption}
\usepackage{subcaption}
\usepackage{tabularx}
\usepackage{stmaryrd}       % \llbracket, \rrbracket
\usepackage{bbm}            % \mathbbm{1}

% ── Page geometry ───────────────────────────────────────────────────────
\usepackage[margin=1in]{geometry}

% ── Theorem environments ────────────────────────────────────────────────
\theoremstyle{plain}
\newtheorem{theorem}{Theorem}[section]
\newtheorem{lemma}[theorem]{Lemma}
\newtheorem{proposition}[theorem]{Proposition}
\newtheorem{corollary}[theorem]{Corollary}

\theoremstyle{definition}
\newtheorem{definition}[theorem]{Definition}
\newtheorem{example}[theorem]{Example}
\newtheorem{construction}[theorem]{Construction}

\theoremstyle{remark}
\newtheorem{remark}[theorem]{Remark}
\newtheorem{notation}[theorem]{Notation}

% ── Python listings style ──────────────────────────────────────────────
\definecolor{jugeoBlue}{HTML}{2563EB}
\definecolor{jugeoGreen}{HTML}{059669}
\definecolor{jugeoRed}{HTML}{DC2626}
\definecolor{jugeoPurple}{HTML}{7C3AED}
\definecolor{jugeoGray}{HTML}{6B7280}
\definecolor{jugeoBg}{HTML}{F8FAFC}

\lstdefinestyle{jugeo-python}{
  language=Python,
  basicstyle=\ttfamily\small,
  keywordstyle=\color{jugeoBlue}\bfseries,
  stringstyle=\color{jugeoGreen},
  commentstyle=\color{jugeoGray}\itshape,
  numberstyle=\tiny\color{jugeoGray},
  backgroundcolor=\color{jugeoBg},
  numbers=left,
  numbersep=6pt,
  frame=single,
  framerule=0.4pt,
  rulecolor=\color{jugeoGray!40},
  breaklines=true,
  breakatwhitespace=true,
  showstringspaces=false,
  tabsize=4,
  xleftmargin=1.5em,
  framexleftmargin=1.5em,
  morekeywords={dataclass,frozen,field,Enum,IntEnum,auto,Any,
                Mapping,Sequence,Optional,match,case},
  literate={→}{{$\to$}}1 {←}{{$\leftarrow$}}1
           {≤}{{$\leq$}}1 {≥}{{$\geq$}}1 {≠}{{$\neq$}}1
           {∈}{{$\in$}}1 {∀}{{$\forall$}}1 {∃}{{$\exists$}}1
           {⊕}{{$\oplus$}}1 {⊖}{{$\ominus$}}1,
  escapeinside={(*@}{@*)},
}
\lstset{style=jugeo-python}

\lstdefinestyle{jugeo-output}{
  basicstyle=\ttfamily\small,
  backgroundcolor=\color{jugeoBg},
  frame=single,
  framerule=0.4pt,
  rulecolor=\color{jugeoGray!40},
  breaklines=true,
  showstringspaces=false,
  xleftmargin=1.5em,
  framexleftmargin=0.5em,
  numbers=none,
}

% ── JuGeo-specific macros ──────────────────────────────────────────────

% Judgment 8-tuple
\newcommand{\J}{\mathcal{J}}
\newcommand{\Jud}[8]{(#1,\,#2,\,#3,\,#4,\,#5,\,#6,\,#7,\,#8)}
\newcommand{\JudShort}[1]{\J_{#1}}

% Components
\newcommand{\coord}{\mathsf{c}}            % coordinate
\newcommand{\prop}{\varphi}                 % proposition
\newcommand{\carrier}{\mathcal{A}}          % carrier type
\newcommand{\evid}{\mathcal{E}}             % evidence bundle
\newcommand{\oblig}{\mathcal{O}}            % obligations
\newcommand{\obstr}{\mathcal{B}}            % obstructions
\newcommand{\trust}{\mathcal{T}}            % trust annotation
\newcommand{\prov}{\Pi}                     % provenance

% Trust algebra
\newcommand{\Talg}{\mathfrak{T}}            % trust ordered algebra
\newcommand{\Eadm}{\mathcal{E}_{\mathrm{adm}}}
\newcommand{\tleq}{\preceq}                 % trust ordering
\newcommand{\tjoin}{\oplus}                 % conservative join
\newcommand{\tatten}{\ominus}               % attenuation
\newcommand{\tpromote}{\uparrow_\pi}        % promotion
\newcommand{\tdemote}{\downarrow_\chi}       % demotion

% Trust tiers
\newcommand{\Tcontra}{\ensuremath{\mathsf{CONTRADICTED}}}
\newcommand{\Tunver}{\ensuremath{\mathsf{UNVERIFIED}}}
\newcommand{\Tcopilot}{\ensuremath{\mathsf{COPILOT}}}
\newcommand{\Toracle}{\ensuremath{\mathsf{ORACLE}}}
\newcommand{\Truntime}{\ensuremath{\mathsf{RUNTIME}}}
\newcommand{\Tsolver}{\ensuremath{\mathsf{SOLVER}}}
\newcommand{\Tproof}{\ensuremath{\mathsf{PROOF}}}

% Site and geometry
\newcommand{\Site}{\mathcal{S}}             % semantic site
\newcommand{\Cov}{\mathrm{Cov}}            % covering family
\newcommand{\Ob}{\mathrm{Ob}}              % objects
\newcommand{\Mor}{\mathrm{Mor}}            % morphisms
\newcommand{\res}{\mathrm{res}}            % restriction
\newcommand{\Cech}{\check{C}}              % Čech complex
\newcommand{\Hcoh}[1]{H^{#1}}             % cohomology group

% Descent
\newcommand{\descent}{\mathrm{desc}}
\newcommand{\glue}{\mathrm{glue}}
\newcommand{\overlap}[2]{#1 \cap #2}

% Semantic moves
\newcommand{\VERIFY}{\textsc{Verify}}
\newcommand{\CONSTRUCT}{\textsc{Construct}}
\newcommand{\NORMALIZE}{\textsc{Normalize}}
\newcommand{\GLUE}{\textsc{Glue}}
\newcommand{\RESTRICT}{\textsc{Restrict}}
\newcommand{\TRANSPORT}{\textsc{Transport}}
\newcommand{\REPAIR}{\textsc{Repair}}
\newcommand{\PROMOTE}{\textsc{Promote}}
\newcommand{\CHALLENGE}{\textsc{Challenge}}
\newcommand{\ESCALATE}{\textsc{Escalate}}

% Fragments
\newcommand{\QFLIA}{\texttt{QF\_LIA}}
\newcommand{\QFLRA}{\texttt{QF\_LRA}}
\newcommand{\QFBV}{\texttt{QF\_BV}}
\newcommand{\QFUF}{\texttt{QF\_UF}}

% Misc
\newcommand{\Python}{\textsc{Python}}
\newcommand{\JuGeo}{\textsc{JuGeo}}
\newcommand{\LEAN}{\textsc{Lean}\,4}
\newcommand{\Fstar}{F$^{\star}$}
\newcommand{\Coq}{\textsc{Coq}}
\newcommand{\Dafny}{\textsc{Dafny}}

% Common shortcuts
\newcommand{\ie}{i.e.\@\xspace}
\newcommand{\eg}{e.g.\@\xspace}
\newcommand{\cf}{cf.\@\xspace}
\newcommand{\wrt}{w.r.t.\@\xspace}

% Evaluation shortcuts
\newcommand{\numcases}{300}
\newcommand{\numspec}{100}
\newcommand{\numeq}{100}
\newcommand{\numbug}{100}
\newcommand{\medtime}{6.5\,ms}
\newcommand{\totaltime}{1.949\,s}


% ── Packages ────────────────────────────────────────────────────────────
\usepackage{amsmath,amssymb,amsthm,mathtools}
\usepackage{tikz-cd}
\usepackage{listings}
\usepackage{xcolor}
\usepackage{xspace}
\usepackage{booktabs}
\usepackage{multirow}
\usepackage{enumitem}
\usepackage[expansion=false]{microtype}
\usepackage{hyperref}
\usepackage{cleveref}
% \usepackage{thmtools}  % disabled: not available in this TeX installation
\usepackage{float}
\usepackage{caption}
\usepackage{subcaption}
\usepackage{tabularx}
\usepackage{stmaryrd}       % \llbracket, \rrbracket
\usepackage{bbm}            % \mathbbm{1}

% ── Page geometry ───────────────────────────────────────────────────────
\usepackage[margin=1in]{geometry}

% ── Theorem environments ────────────────────────────────────────────────
\theoremstyle{plain}
\newtheorem{theorem}{Theorem}[section]
\newtheorem{lemma}[theorem]{Lemma}
\newtheorem{proposition}[theorem]{Proposition}
\newtheorem{corollary}[theorem]{Corollary}

\theoremstyle{definition}
\newtheorem{definition}[theorem]{Definition}
\newtheorem{example}[theorem]{Example}
\newtheorem{construction}[theorem]{Construction}

\theoremstyle{remark}
\newtheorem{remark}[theorem]{Remark}
\newtheorem{notation}[theorem]{Notation}

% ── Python listings style ──────────────────────────────────────────────
\definecolor{jugeoBlue}{HTML}{2563EB}
\definecolor{jugeoGreen}{HTML}{059669}
\definecolor{jugeoRed}{HTML}{DC2626}
\definecolor{jugeoPurple}{HTML}{7C3AED}
\definecolor{jugeoGray}{HTML}{6B7280}
\definecolor{jugeoBg}{HTML}{F8FAFC}

\lstdefinestyle{jugeo-python}{
  language=Python,
  basicstyle=\ttfamily\small,
  keywordstyle=\color{jugeoBlue}\bfseries,
  stringstyle=\color{jugeoGreen},
  commentstyle=\color{jugeoGray}\itshape,
  numberstyle=\tiny\color{jugeoGray},
  backgroundcolor=\color{jugeoBg},
  numbers=left,
  numbersep=6pt,
  frame=single,
  framerule=0.4pt,
  rulecolor=\color{jugeoGray!40},
  breaklines=true,
  breakatwhitespace=true,
  showstringspaces=false,
  tabsize=4,
  xleftmargin=1.5em,
  framexleftmargin=1.5em,
  morekeywords={dataclass,frozen,field,Enum,IntEnum,auto,Any,
                Mapping,Sequence,Optional,match,case},
  literate={→}{{$\to$}}1 {←}{{$\leftarrow$}}1
           {≤}{{$\leq$}}1 {≥}{{$\geq$}}1 {≠}{{$\neq$}}1
           {∈}{{$\in$}}1 {∀}{{$\forall$}}1 {∃}{{$\exists$}}1
           {⊕}{{$\oplus$}}1 {⊖}{{$\ominus$}}1,
  escapeinside={(*@}{@*)},
}
\lstset{style=jugeo-python}

\lstdefinestyle{jugeo-output}{
  basicstyle=\ttfamily\small,
  backgroundcolor=\color{jugeoBg},
  frame=single,
  framerule=0.4pt,
  rulecolor=\color{jugeoGray!40},
  breaklines=true,
  showstringspaces=false,
  xleftmargin=1.5em,
  framexleftmargin=0.5em,
  numbers=none,
}

% ── JuGeo-specific macros ──────────────────────────────────────────────

% Judgment 8-tuple
\newcommand{\J}{\mathcal{J}}
\newcommand{\Jud}[8]{(#1,\,#2,\,#3,\,#4,\,#5,\,#6,\,#7,\,#8)}
\newcommand{\JudShort}[1]{\J_{#1}}

% Components
\newcommand{\coord}{\mathsf{c}}            % coordinate
\newcommand{\prop}{\varphi}                 % proposition
\newcommand{\carrier}{\mathcal{A}}          % carrier type
\newcommand{\evid}{\mathcal{E}}             % evidence bundle
\newcommand{\oblig}{\mathcal{O}}            % obligations
\newcommand{\obstr}{\mathcal{B}}            % obstructions
\newcommand{\trust}{\mathcal{T}}            % trust annotation
\newcommand{\prov}{\Pi}                     % provenance

% Trust algebra
\newcommand{\Talg}{\mathfrak{T}}            % trust ordered algebra
\newcommand{\Eadm}{\mathcal{E}_{\mathrm{adm}}}
\newcommand{\tleq}{\preceq}                 % trust ordering
\newcommand{\tjoin}{\oplus}                 % conservative join
\newcommand{\tatten}{\ominus}               % attenuation
\newcommand{\tpromote}{\uparrow_\pi}        % promotion
\newcommand{\tdemote}{\downarrow_\chi}       % demotion

% Trust tiers
\newcommand{\Tcontra}{\ensuremath{\mathsf{CONTRADICTED}}}
\newcommand{\Tunver}{\ensuremath{\mathsf{UNVERIFIED}}}
\newcommand{\Tcopilot}{\ensuremath{\mathsf{COPILOT}}}
\newcommand{\Toracle}{\ensuremath{\mathsf{ORACLE}}}
\newcommand{\Truntime}{\ensuremath{\mathsf{RUNTIME}}}
\newcommand{\Tsolver}{\ensuremath{\mathsf{SOLVER}}}
\newcommand{\Tproof}{\ensuremath{\mathsf{PROOF}}}

% Site and geometry
\newcommand{\Site}{\mathcal{S}}             % semantic site
\newcommand{\Cov}{\mathrm{Cov}}            % covering family
\newcommand{\Ob}{\mathrm{Ob}}              % objects
\newcommand{\Mor}{\mathrm{Mor}}            % morphisms
\newcommand{\res}{\mathrm{res}}            % restriction
\newcommand{\Cech}{\check{C}}              % Čech complex
\newcommand{\Hcoh}[1]{H^{#1}}             % cohomology group

% Descent
\newcommand{\descent}{\mathrm{desc}}
\newcommand{\glue}{\mathrm{glue}}
\newcommand{\overlap}[2]{#1 \cap #2}

% Semantic moves
\newcommand{\VERIFY}{\textsc{Verify}}
\newcommand{\CONSTRUCT}{\textsc{Construct}}
\newcommand{\NORMALIZE}{\textsc{Normalize}}
\newcommand{\GLUE}{\textsc{Glue}}
\newcommand{\RESTRICT}{\textsc{Restrict}}
\newcommand{\TRANSPORT}{\textsc{Transport}}
\newcommand{\REPAIR}{\textsc{Repair}}
\newcommand{\PROMOTE}{\textsc{Promote}}
\newcommand{\CHALLENGE}{\textsc{Challenge}}
\newcommand{\ESCALATE}{\textsc{Escalate}}

% Fragments
\newcommand{\QFLIA}{\texttt{QF\_LIA}}
\newcommand{\QFLRA}{\texttt{QF\_LRA}}
\newcommand{\QFBV}{\texttt{QF\_BV}}
\newcommand{\QFUF}{\texttt{QF\_UF}}

% Misc
\newcommand{\Python}{\textsc{Python}}
\newcommand{\JuGeo}{\textsc{JuGeo}}
\newcommand{\LEAN}{\textsc{Lean}\,4}
\newcommand{\Fstar}{F$^{\star}$}
\newcommand{\Coq}{\textsc{Coq}}
\newcommand{\Dafny}{\textsc{Dafny}}

% Common shortcuts
\newcommand{\ie}{i.e.\@\xspace}
\newcommand{\eg}{e.g.\@\xspace}
\newcommand{\cf}{cf.\@\xspace}
\newcommand{\wrt}{w.r.t.\@\xspace}

% Evaluation shortcuts
\newcommand{\numcases}{300}
\newcommand{\numspec}{100}
\newcommand{\numeq}{100}
\newcommand{\numbug}{100}
\newcommand{\medtime}{6.5\,ms}
\newcommand{\totaltime}{1.949\,s}


% ── Packages ────────────────────────────────────────────────────────────
\usepackage{amsmath,amssymb,amsthm,mathtools}
\usepackage{tikz-cd}
\usepackage{listings}
\usepackage{xcolor}
\usepackage{xspace}
\usepackage{booktabs}
\usepackage{multirow}
\usepackage{enumitem}
\usepackage[expansion=false]{microtype}
\usepackage{hyperref}
\usepackage{cleveref}
% \usepackage{thmtools}  % disabled: not available in this TeX installation
\usepackage{float}
\usepackage{caption}
\usepackage{subcaption}
\usepackage{tabularx}
\usepackage{stmaryrd}       % \llbracket, \rrbracket
\usepackage{bbm}            % \mathbbm{1}

% ── Page geometry ───────────────────────────────────────────────────────
\usepackage[margin=1in]{geometry}

% ── Theorem environments ────────────────────────────────────────────────
\theoremstyle{plain}
\newtheorem{theorem}{Theorem}[section]
\newtheorem{lemma}[theorem]{Lemma}
\newtheorem{proposition}[theorem]{Proposition}
\newtheorem{corollary}[theorem]{Corollary}

\theoremstyle{definition}
\newtheorem{definition}[theorem]{Definition}
\newtheorem{example}[theorem]{Example}
\newtheorem{construction}[theorem]{Construction}

\theoremstyle{remark}
\newtheorem{remark}[theorem]{Remark}
\newtheorem{notation}[theorem]{Notation}

% ── Python listings style ──────────────────────────────────────────────
\definecolor{jugeoBlue}{HTML}{2563EB}
\definecolor{jugeoGreen}{HTML}{059669}
\definecolor{jugeoRed}{HTML}{DC2626}
\definecolor{jugeoPurple}{HTML}{7C3AED}
\definecolor{jugeoGray}{HTML}{6B7280}
\definecolor{jugeoBg}{HTML}{F8FAFC}

\lstdefinestyle{jugeo-python}{
  language=Python,
  basicstyle=\ttfamily\small,
  keywordstyle=\color{jugeoBlue}\bfseries,
  stringstyle=\color{jugeoGreen},
  commentstyle=\color{jugeoGray}\itshape,
  numberstyle=\tiny\color{jugeoGray},
  backgroundcolor=\color{jugeoBg},
  numbers=left,
  numbersep=6pt,
  frame=single,
  framerule=0.4pt,
  rulecolor=\color{jugeoGray!40},
  breaklines=true,
  breakatwhitespace=true,
  showstringspaces=false,
  tabsize=4,
  xleftmargin=1.5em,
  framexleftmargin=1.5em,
  morekeywords={dataclass,frozen,field,Enum,IntEnum,auto,Any,
                Mapping,Sequence,Optional,match,case},
  literate={→}{{$\to$}}1 {←}{{$\leftarrow$}}1
           {≤}{{$\leq$}}1 {≥}{{$\geq$}}1 {≠}{{$\neq$}}1
           {∈}{{$\in$}}1 {∀}{{$\forall$}}1 {∃}{{$\exists$}}1
           {⊕}{{$\oplus$}}1 {⊖}{{$\ominus$}}1,
  escapeinside={(*@}{@*)},
}
\lstset{style=jugeo-python}

\lstdefinestyle{jugeo-output}{
  basicstyle=\ttfamily\small,
  backgroundcolor=\color{jugeoBg},
  frame=single,
  framerule=0.4pt,
  rulecolor=\color{jugeoGray!40},
  breaklines=true,
  showstringspaces=false,
  xleftmargin=1.5em,
  framexleftmargin=0.5em,
  numbers=none,
}

% ── JuGeo-specific macros ──────────────────────────────────────────────

% Judgment 8-tuple
\newcommand{\J}{\mathcal{J}}
\newcommand{\Jud}[8]{(#1,\,#2,\,#3,\,#4,\,#5,\,#6,\,#7,\,#8)}
\newcommand{\JudShort}[1]{\J_{#1}}

% Components
\newcommand{\coord}{\mathsf{c}}            % coordinate
\newcommand{\prop}{\varphi}                 % proposition
\newcommand{\carrier}{\mathcal{A}}          % carrier type
\newcommand{\evid}{\mathcal{E}}             % evidence bundle
\newcommand{\oblig}{\mathcal{O}}            % obligations
\newcommand{\obstr}{\mathcal{B}}            % obstructions
\newcommand{\trust}{\mathcal{T}}            % trust annotation
\newcommand{\prov}{\Pi}                     % provenance

% Trust algebra
\newcommand{\Talg}{\mathfrak{T}}            % trust ordered algebra
\newcommand{\Eadm}{\mathcal{E}_{\mathrm{adm}}}
\newcommand{\tleq}{\preceq}                 % trust ordering
\newcommand{\tjoin}{\oplus}                 % conservative join
\newcommand{\tatten}{\ominus}               % attenuation
\newcommand{\tpromote}{\uparrow_\pi}        % promotion
\newcommand{\tdemote}{\downarrow_\chi}       % demotion

% Trust tiers
\newcommand{\Tcontra}{\ensuremath{\mathsf{CONTRADICTED}}}
\newcommand{\Tunver}{\ensuremath{\mathsf{UNVERIFIED}}}
\newcommand{\Tcopilot}{\ensuremath{\mathsf{COPILOT}}}
\newcommand{\Toracle}{\ensuremath{\mathsf{ORACLE}}}
\newcommand{\Truntime}{\ensuremath{\mathsf{RUNTIME}}}
\newcommand{\Tsolver}{\ensuremath{\mathsf{SOLVER}}}
\newcommand{\Tproof}{\ensuremath{\mathsf{PROOF}}}

% Site and geometry
\newcommand{\Site}{\mathcal{S}}             % semantic site
\newcommand{\Cov}{\mathrm{Cov}}            % covering family
\newcommand{\Ob}{\mathrm{Ob}}              % objects
\newcommand{\Mor}{\mathrm{Mor}}            % morphisms
\newcommand{\res}{\mathrm{res}}            % restriction
\newcommand{\Cech}{\check{C}}              % Čech complex
\newcommand{\Hcoh}[1]{H^{#1}}             % cohomology group

% Descent
\newcommand{\descent}{\mathrm{desc}}
\newcommand{\glue}{\mathrm{glue}}
\newcommand{\overlap}[2]{#1 \cap #2}

% Semantic moves
\newcommand{\VERIFY}{\textsc{Verify}}
\newcommand{\CONSTRUCT}{\textsc{Construct}}
\newcommand{\NORMALIZE}{\textsc{Normalize}}
\newcommand{\GLUE}{\textsc{Glue}}
\newcommand{\RESTRICT}{\textsc{Restrict}}
\newcommand{\TRANSPORT}{\textsc{Transport}}
\newcommand{\REPAIR}{\textsc{Repair}}
\newcommand{\PROMOTE}{\textsc{Promote}}
\newcommand{\CHALLENGE}{\textsc{Challenge}}
\newcommand{\ESCALATE}{\textsc{Escalate}}

% Fragments
\newcommand{\QFLIA}{\texttt{QF\_LIA}}
\newcommand{\QFLRA}{\texttt{QF\_LRA}}
\newcommand{\QFBV}{\texttt{QF\_BV}}
\newcommand{\QFUF}{\texttt{QF\_UF}}

% Misc
\newcommand{\Python}{\textsc{Python}}
\newcommand{\JuGeo}{\textsc{JuGeo}}
\newcommand{\LEAN}{\textsc{Lean}\,4}
\newcommand{\Fstar}{F$^{\star}$}
\newcommand{\Coq}{\textsc{Coq}}
\newcommand{\Dafny}{\textsc{Dafny}}

% Common shortcuts
\newcommand{\ie}{i.e.\@\xspace}
\newcommand{\eg}{e.g.\@\xspace}
\newcommand{\cf}{cf.\@\xspace}
\newcommand{\wrt}{w.r.t.\@\xspace}

% Evaluation shortcuts
\newcommand{\numcases}{300}
\newcommand{\numspec}{100}
\newcommand{\numeq}{100}
\newcommand{\numbug}{100}
\newcommand{\medtime}{6.5\,ms}
\newcommand{\totaltime}{1.949\,s}

% experiment-data.tex — AUTO-GENERATED from real jugeo CLI runs
% DO NOT EDIT — regenerate with: python3 experiments/run_universal_metrics.py
% Generated from 30 programs across 6 domains

\newcommand{\expTotalPrograms}{30}
\newcommand{\expVerified}{30}
\newcommand{\expAccuracy}{100.0\%}
\newcommand{\expCoordsMin}{2}
\newcommand{\expCoordsMax}{10}
\newcommand{\expCoordsMean}{5.2}
\newcommand{\expCoordsMedian}{5.5}
\newcommand{\expPropsMin}{4}
\newcommand{\expPropsMax}{20}
\newcommand{\expPropsMean}{10.4}
\newcommand{\expPropsSum}{311}
\newcommand{\expPropsOkSum}{310}
\newcommand{\expTimeMin}{3.506\,s}
\newcommand{\expTimeMax}{6.714\,s}
\newcommand{\expTimeMean}{4.64\,s}
\newcommand{\expTimeMedian}{4.508\,s}
\newcommand{\expTimeTotal}{139.204\,s}
\newcommand{\expObstructionTotal}{0}
\newcommand{\expHOneZero}{30}
\newcommand{\expDomainCount}{6}
\newcommand{\expTrustCOPILOTSUGGESTED}{30}
\newcommand{\expDomainDsCount}{5}
\newcommand{\expDomainDsVerified}{5}
\newcommand{\expDomainDsAvgCoords}{7.6}
\newcommand{\expDomainDsAvgProps}{15.2}
\newcommand{\expDomainMathCount}{5}
\newcommand{\expDomainMathVerified}{5}
\newcommand{\expDomainMathAvgCoords}{4.8}
\newcommand{\expDomainMathAvgProps}{9.4}
\newcommand{\expDomainSmCount}{5}
\newcommand{\expDomainSmVerified}{5}
\newcommand{\expDomainSmAvgCoords}{7.6}
\newcommand{\expDomainSmAvgProps}{15.2}
\newcommand{\expDomainSortCount}{5}
\newcommand{\expDomainSortVerified}{5}
\newcommand{\expDomainSortAvgCoords}{2.4}
\newcommand{\expDomainSortAvgProps}{4.8}
\newcommand{\expDomainStrCount}{5}
\newcommand{\expDomainStrVerified}{5}
\newcommand{\expDomainStrAvgCoords}{3.2}
\newcommand{\expDomainStrAvgProps}{6.4}
\newcommand{\expDomainWebCount}{5}
\newcommand{\expDomainWebVerified}{5}
\newcommand{\expDomainWebAvgCoords}{5.6}
\newcommand{\expDomainWebAvgProps}{11.2}

% subsystem-data.tex — AUTO-GENERATED from real jugeo subsystem experiments
% DO NOT EDIT — regenerate with: python3 experiments/run_subsystem_experiments.py

\newcommand{\subCliTotal}{10}
\newcommand{\subCliVerified}{9}
\newcommand{\subCliAccuracy}{90.0\%}
\newcommand{\subCliMeanTime}{4.132\,s}
\newcommand{\subCliMinTime}{3.472\,s}
\newcommand{\subCliMaxTime}{5.372\,s}
\newcommand{\subCliTotalCoords}{54}
\newcommand{\subCliTotalProps}{107}
\newcommand{\subCliTotalPropsOk}{106}
\newcommand{\subSiteCalls}{90}
\newcommand{\subSiteSuccessful}{69}
\newcommand{\subSiteSuccessRate}{76.7\%}
\newcommand{\subSiteMeanTime}{0.0001\,s}
\newcommand{\subSiteTotalTime}{0.005\,s}
\newcommand{\subSpecificationsatisfactionSmallTime}{0.0003\,s}
\newcommand{\subSpecificationsatisfactionMediumTime}{0.0001\,s}
\newcommand{\subSpecificationsatisfactionLargeTime}{0.0001\,s}
\newcommand{\subInhabitantfleetSmallTime}{0.0\,s}
\newcommand{\subInhabitantfleetMediumTime}{0.0\,s}
\newcommand{\subInhabitantfleetLargeTime}{0.0\,s}
\newcommand{\subTheoremecologySmallTime}{0.0\,s}
\newcommand{\subTheoremecologyMediumTime}{0.0\,s}
\newcommand{\subTheoremecologyLargeTime}{0.0\,s}
\newcommand{\subStatespaceexplorationSmallTime}{0.0\,s}
\newcommand{\subStatespaceexplorationMediumTime}{0.0\,s}
\newcommand{\subStatespaceexplorationLargeTime}{0.0\,s}
\newcommand{\subRepairsemanticsSmallTime}{0.0\,s}
\newcommand{\subRepairsemanticsMediumTime}{0.0\,s}
\newcommand{\subRepairsemanticsLargeTime}{0.0\,s}
\newcommand{\subReplaygluingSmallTime}{0.0\,s}
\newcommand{\subReplaygluingMediumTime}{0.0\,s}
\newcommand{\subReplaygluingLargeTime}{0.0\,s}
\newcommand{\subSemanticfuturesSmallTime}{0.0\,s}
\newcommand{\subSemanticfuturesMediumTime}{0.0\,s}
\newcommand{\subSemanticfuturesLargeTime}{0.0\,s}
\newcommand{\subHypercovertreatySmallTime}{0.0\,s}
\newcommand{\subHypercovertreatyMediumTime}{0.0\,s}
\newcommand{\subHypercovertreatyLargeTime}{0.0\,s}
\newcommand{\subEncodeforsolverSmallTime}{0.0026\,s}
\newcommand{\subEncodeforsolverMediumTime}{0.0002\,s}
\newcommand{\subEncodeforsolverLargeTime}{0.0001\,s}
\newcommand{\subInterfaceroutingSmallTime}{0.0\,s}
\newcommand{\subInterfaceroutingMediumTime}{0.0\,s}
\newcommand{\subInterfaceroutingLargeTime}{0.0\,s}
\newcommand{\subOrchestrateverificationSmallTime}{0.0002\,s}
\newcommand{\subOrchestrateverificationMediumTime}{0.0\,s}
\newcommand{\subOrchestrateverificationLargeTime}{0.0\,s}
\newcommand{\subRunfulldescentSmallTime}{0.0\,s}
\newcommand{\subRunfulldescentMediumTime}{0.0\,s}
\newcommand{\subRunfulldescentLargeTime}{0.0\,s}
\newcommand{\subKernellifecycleSmallTime}{0.0\,s}
\newcommand{\subKernellifecycleMediumTime}{0.0\,s}
\newcommand{\subKernellifecycleLargeTime}{0.0\,s}
\newcommand{\subDiscoverypipelineSmallTime}{0.0\,s}
\newcommand{\subDiscoverypipelineMediumTime}{0.0\,s}
\newcommand{\subDiscoverypipelineLargeTime}{0.0\,s}
\newcommand{\subAnalogytransportSmallTime}{0.0\,s}
\newcommand{\subAnalogytransportMediumTime}{0.0\,s}
\newcommand{\subAnalogytransportLargeTime}{0.0\,s}
\newcommand{\subFormalcoresiteSmallTime}{0.0001\,s}
\newcommand{\subFormalcoresiteMediumTime}{0.0\,s}
\newcommand{\subFormalcoresiteLargeTime}{0.0\,s}
\newcommand{\subProblematlasSmallTime}{0.0\,s}
\newcommand{\subProblematlasMediumTime}{0.0\,s}
\newcommand{\subProblematlasLargeTime}{0.0\,s}
\newcommand{\subPublicalignmentSmallTime}{0.0\,s}
\newcommand{\subPublicalignmentMediumTime}{0.0\,s}
\newcommand{\subPublicalignmentLargeTime}{0.0\,s}
\newcommand{\subRegimebootstrappingSmallTime}{0.0\,s}
\newcommand{\subRegimebootstrappingMediumTime}{0.0\,s}
\newcommand{\subRegimebootstrappingLargeTime}{0.0\,s}
\newcommand{\subRelationalrefinementSmallTime}{0.0\,s}
\newcommand{\subRelationalrefinementMediumTime}{0.0\,s}
\newcommand{\subRelationalrefinementLargeTime}{0.0\,s}
\newcommand{\subSemanticclosureSmallTime}{0.0\,s}
\newcommand{\subSemanticclosureMediumTime}{0.0\,s}
\newcommand{\subSemanticclosureLargeTime}{0.0\,s}
\newcommand{\subTheoremeconomicsSmallTime}{0.0001\,s}
\newcommand{\subTheoremeconomicsMediumTime}{0.0\,s}
\newcommand{\subTheoremeconomicsLargeTime}{0.0\,s}
\newcommand{\subBenchmarksuiteSmallTime}{0.0\,s}
\newcommand{\subBenchmarksuiteMediumTime}{0.0\,s}
\newcommand{\subBenchmarksuiteLargeTime}{0.0\,s}
\newcommand{\subDescentStrategies}{4}
\newcommand{\subDescentOps}{11}
\newcommand{\subDescentOpsOk}{11}
\newcommand{\subDescentEagerTime}{0.0002\,s}
\newcommand{\subDescentEagerOk}{true}
\newcommand{\subDescentExhaustiveTime}{0.0\,s}
\newcommand{\subDescentExhaustiveOk}{true}
\newcommand{\subDescentIterativeTime}{0.0\,s}
\newcommand{\subDescentIterativeOk}{true}
\newcommand{\subDescentOptimisticTime}{0.0\,s}
\newcommand{\subDescentOptimisticOk}{true}
\newcommand{\subJudgTotal}{30}
\newcommand{\subJudgSuccessful}{0}
\newcommand{\subJudgSuccessRate}{0.0\%}
\newcommand{\subJudgMeanTimeUs}{7.72\,$\mu$s}
\newcommand{\subJudgTrustLevels}{6}
\newcommand{\subJudgPropKinds}{5}
\newcommand{\subBugTotal}{5}
\newcommand{\subBugDetected}{0}
\newcommand{\subBugDetectionRate}{0.0\%}
\newcommand{\subBugFalsePositiveRate}{0.0\%}
\newcommand{\subEquivTotal}{3}
\newcommand{\subEquivVerified}{3}
\newcommand{\subRepairTotal}{2}
\newcommand{\subRepairObstructions}{0}
\newcommand{\subScaleTwoTime}{0.0001\,s}
\newcommand{\subScaleFourTime}{0.0\,s}
\newcommand{\subScaleEightTime}{0.0\,s}
\newcommand{\subScaleTwelveTime}{0.0\,s}
\newcommand{\subScaleSixteenTime}{0.0\,s}
\newcommand{\subScaleTwentyTime}{0.0\,s}
\newcommand{\subTotalExperimentTime}{95.6\,s}


%% -----------------------------------------------------------------------
%% Paper 33 — Text and String Encodings: SMT Theories for Python String
%% Verification
%% Paper-specific macros (none from jugeo-common redefined)
%% -----------------------------------------------------------------------
\newcommand{\StrEnc}{\ensuremath{\mathsf{StrEnc}}}
\newcommand{\SymStr}{\ensuremath{\mathsf{SymStr}}}
\newcommand{\NamLaw}{\ensuremath{\mathsf{NamLaw}}}
\newcommand{\DocShad}{\ensuremath{\mathsf{DocShad}}}
\newcommand{\StreamEnc}{\ensuremath{\mathsf{StreamEnc}}}
\newcommand{\TextEnc}{\ensuremath{\mathsf{TextEnc}}}
\newcommand{\QFSLIA}{\textsf{QF\_SLIA}}
\newcommand{\QFS}{\textsf{QF\_S}}
\newcommand{\strsort}{\ensuremath{\mathtt{String}}}
\newcommand{\strlen}[1]{\ensuremath{\mathsf{str.len}(#1)}}
\newcommand{\strpre}[2]{\ensuremath{\mathsf{str.prefixof}(#1,\,#2)}}
\newcommand{\strsuf}[2]{\ensuremath{\mathsf{str.suffixof}(#1,\,#2)}}
\newcommand{\strinre}[2]{\ensuremath{\mathsf{str.in\_re}(#1,\,#2)}}
\newcommand{\strconcat}[2]{\ensuremath{#1 \mathbin{+\!\!+} #2}}
\newcommand{\strcont}[2]{\ensuremath{\mathsf{str.contains}(#1,\,#2)}}
\newcommand{\NFC}{\ensuremath{\mathrm{NFC}}}
\newcommand{\NFKC}{\ensuremath{\mathrm{NFKC}}}
\newcommand{\StrOp}{\ensuremath{\mathsf{StrOp}}}
\newcommand{\CKind}{\ensuremath{\mathsf{CKind}}}
\newcommand{\CStr}{\ensuremath{\mathsf{CStr}}}
\newcommand{\PropPair}[2]{\ensuremath{\langle #1,\,#2 \rangle}}
\newcommand{\JGProp}{\ensuremath{\mathsf{Prop}}}
\newcommand{\Chunk}{\ensuremath{\mathsf{Chunk}}}
\newcommand{\strreplace}[3]{\ensuremath{\mathsf{str.replace}(#1,\,#2,\,#3)}}
\newcommand{\strindexof}[2]{\ensuremath{\mathsf{str.indexof}(#1,\,#2)}}
\newcommand{\strsubstr}[3]{\ensuremath{\mathsf{str.substr}(#1,\,#2,\,#3)}}
\newcommand{\restar}[1]{\ensuremath{#1^{*}}}
\newcommand{\replus}[1]{\ensuremath{#1^{+}}}
\newcommand{\reopt}[1]{\ensuremath{#1^{?}}}
\newcommand{\rerange}[2]{\ensuremath{\mathsf{re.range}(#1,\,#2)}}
\newcommand{\reunion}[2]{\ensuremath{#1 \cup #2}}
\newcommand{\recat}[2]{\ensuremath{#1 \cdot #2}}
\newcommand{\encode}{\ensuremath{\mathsf{enc}}}
\newcommand{\bigcat}{\ensuremath{\bigoplus}}

\title{\textbf{Text and String Encodings:\\
       SMT Theories for Python String Verification}}
\author{Halley Young \\
        \textit{Microsoft Research} \\
        \texttt{halleyyoung@microsoft.com}}
\date{\today}

\begin{document}
\maketitle

\begin{abstract}
String manipulation is pervasive in Python programs: identifiers follow naming
conventions, format strings carry structural constraints, documentation embeds
preconditions, and streaming pipelines process text in chunks.  Yet
string-heavy programs are notoriously difficult to verify using arithmetic-only
solvers.  We present \textbf{\TextEnc}, the string encoding layer of the
\JuGeo{} verification framework, which translates Python \texttt{str}
operations into the \QFSLIA{} SMT theory understood by Z3.  The layer
comprises three cooperating components: (i)~a \emph{string model}
(\StrEnc{}/\SymStr) that captures Python's \texttt{str} semantics including
Unicode normalisation, (ii)~a \emph{streaming encoding} (\StreamEnc) that
handles string processing in chunks, and (iii)~a \emph{family classifier} that
routes each constraint to its least complex decidable fragment.  We state and
prove a \emph{String Encoding Completeness Theorem}: every Python string
operation has a corresponding \QFSLIA{} encoding that preserves equality
semantics.  Experiments on \numcases{} Python functions demonstrate that
\TextEnc{} reduces average query time compared to na\"ive bit-vector
encodings while maintaining full soundness.
\end{abstract}

\tableofcontents
\newpage

%%=======================================================================
\section{Introduction}
\label{sec:intro}
%%=======================================================================

Python's \texttt{str} type is one of the most heavily used data types in the
language.  Programs name variables with identifiers, route HTTP requests via
URL strings, persist data through serialisation formats, and generate user
output via f-strings and \texttt{format} templates.  The semantics of
\texttt{str} operations—concatenation, slicing, regular-expression matching,
Unicode normalisation—are precisely specified in the Python language
reference~\cite{python-ref}, yet most verification tools treat strings as
uninterpreted blobs or fall back to costly bit-vector models.

\paragraph{The string verification gap.}
Arithmetic SMT theories (\QFLIA, \QFLRA) express numeric invariants with
great precision.  Bit-vector theories (\QFBV) model machine-word arithmetic.
The SMT community has more recently standardised the \emph{theory of strings}
(\QFS{}/\QFSLIA)~\cite{smtlib-strings}, which Z3~\cite{z3} and
CVC5~\cite{cvc5} implement.  However, the gap between Python's rich string
semantics and the raw SMT string theory primitives is non-trivial: Python
strings are Unicode sequences (not byte arrays), \texttt{str.find} returns
$-1$ on failure (not an exception), and f-string formatting composes multiple
operations that must be encoded atomically.  Bridging this gap requires a
carefully engineered encoding layer.

\paragraph{Contributions.}
This paper makes the following contributions.
\begin{enumerate}
  \item \textbf{String Model (\S\ref{sec:model})}: We define the \StrEnc{}
    and \SymStr{} data structures that capture Python \texttt{str} semantics
    in \QFSLIA{} (\S\ref{sec:model}).

  \item \textbf{Core Operations (\S\ref{sec:ops})}: We give precise \QFSLIA{}
    encodings for concatenation, slicing, \texttt{find}/\texttt{index},
    \texttt{replace}, and format strings.

  \item \textbf{Streaming Encoding (\S\ref{sec:streaming})}: We introduce
    \StreamEnc{}, which processes strings in chunks while preserving
    whole-string invariants.

  \item \textbf{Unicode Handling (\S\ref{sec:unicode})}: We explain how the
    \texttt{NormalizedTextEnvironment} enforces NFC invariants in the
    encoding.

  \item \textbf{Regex Integration (\S\ref{sec:regex})}: We show how Python
    regex patterns are compiled to \QFSLIA{} regular-expression terms.

  \item \textbf{Completeness Theorem (\S\ref{sec:theorem})}: We prove that
    every Python string operation has a corresponding \QFSLIA{} encoding that
    preserves equality semantics, and provide the Lean~4 formalisation in
    \texttt{Paper33\_TextEncodings.lean}.

  \item \textbf{Evaluation (\S\ref{sec:experiments})}: We measure overhead and
    solver performance on the \JuGeo{} benchmark suite.
\end{enumerate}

\paragraph{Relation to earlier papers.}
Papers~1--5 established the judgment 8-tuple $\Jud{\coord}{\prop}{\carrier}%
{\evid}{\oblig}{\obstr}{\trust}{\prov}$ and trust algebra.  Papers~6--10
gave Python-specific encodings for arithmetic, effects, and proof-carrying
contracts.  Papers~11--32 applied the framework to advanced algebraic and
geometric settings.  The present paper focuses on the \emph{string} component
of $\evid$ and $\oblig$, which previous papers treated only lightly.

%%=======================================================================
\section{String Model}
\label{sec:model}
%%=======================================================================

\subsection{Python String Semantics}
\label{sec:py-str}

Python's \texttt{str} is an immutable sequence of Unicode code points.
Unlike C strings, it carries no null terminator; unlike Java \texttt{String},
it has no UTF-16 surrogate-pair artefacts.  The key semantic properties we
must capture are:

\begin{enumerate}
  \item \textbf{Unicode}: every character is a code point
    $c \in [0, 0x10\mathrm{FFFF}]$.
  \item \textbf{Immutability}: operations return new strings; no in-place
    mutation.
  \item \textbf{Length}: $\strlen{s}$ returns the number of code points, not
    bytes.
  \item \textbf{Slicing}: \texttt{s[i:j]} equals
    $\strsubstr{s}{i}{j - i}$ with clamping semantics for out-of-bounds
    indices.
  \item \textbf{Find}: \texttt{s.find(t)} returns the first index $i$ such
    that $\strsubstr{s}{i}{\strlen{t}} = t$, or $-1$ if absent.
  \item \textbf{Replace}: \texttt{s.replace(old, new)} replaces the first
    (or all, with \texttt{count}) occurrences.
  \item \textbf{Format}: f-strings and \texttt{str.format} compose substrings
    and conversions.
\end{enumerate}

The \QFSLIA{} theory provides sorts $\strsort$ and $\mathtt{Int}$ with mixed
string-integer operators.  The key SMT-LIB2 primitives we use are listed in
Table~\ref{tab:smtprims}.

\begin{table}[h]
\centering
\caption{SMT-LIB2 string primitives used by \TextEnc.}
\label{tab:smtprims}
\begin{tabular}{lll}
\hline
\textbf{SMT Primitive} & \textbf{Python Equivalent} & \textbf{Fragment} \\
\hline
$\strlen{s}$                   & \texttt{len(s)}              & \QFSLIA \\
$\strpre{p}{s}$                & \texttt{s.startswith(p)}     & \QFS \\
$\strsuf{p}{s}$                & \texttt{s.endswith(p)}       & \QFS \\
$\strcont{s}{t}$               & \texttt{t in s}              & \QFS \\
$\strinre{s}{r}$               & \texttt{re.fullmatch(r, s)}  & \QFS \\
$\strconcat{s}{t}$             & \texttt{s + t}               & \QFS \\
$\strsubstr{s}{i}{n}$          & \texttt{s[i:i+n]}            & \QFSLIA \\
$\strindexof{s}{t}$            & \texttt{s.find(t)}           & \QFSLIA \\
$\strreplace{s}{\mathit{old}}{\mathit{new}}$ & \texttt{s.replace(old, new)} & \QFS \\
\hline
\end{tabular}
\end{table}

\subsection{The \texttt{StrEnc} Data Structure}
\label{sec:strenc}

\begin{definition}[String Encoding]
\label{def:strenc}
A \emph{string encoding} is a frozen record
\[
  \StrEnc = \bigl(
    v,\; z,\; \ell^?,\;
    \vec{p},\; \vec{s},\; \vec{r},\; \pi
  \bigr)
\]
where:
\begin{itemize}
  \item $v \in \Sigma^*$ is the concrete string value (may be empty for a
    purely symbolic encoding),
  \item $z$ is the SMT-LIB2 variable name for the string sort,
  \item $\ell^? \in \mathbb{Z}_{\geq 0} \cup \{\bot\}$ is an optional explicit
    length constraint,
  \item $\vec{p} = (p_1, \ldots, p_k)$ are required prefixes,
  \item $\vec{s} = (s_1, \ldots, s_m)$ are required suffixes,
  \item $\vec{r} = (r_1, \ldots, r_n)$ are regular-expression membership
    constraints, and
  \item $\pi$ is a provenance tag for debugging.
\end{itemize}
The \emph{assertion} of $\StrEnc$ is the conjunction
\[
  \phi(\StrEnc) \;:=\;
    \bigl(\ell^? \neq \bot \Rightarrow \strlen{z} = \ell^?\bigr)
    \;\wedge\;
    \bigwedge_i \strpre{p_i}{z}
    \;\wedge\;
    \bigwedge_j \strsuf{s_j}{z}
    \;\wedge\;
    \bigwedge_k \strinre{z}{r_k}.
\]
\end{definition}

\StrEnc{} is implemented as a frozen Python \texttt{dataclass} so that
instances can be shared across solver sessions without accidental mutation.

\subsection{Symbolic Text Variables}
\label{sec:symstr}

When the concrete value is not known at encoding time—for example, a function
parameter whose value is determined by the caller—we use \emph{symbolic text}
variables.

\begin{definition}[Symbolic Text]
\label{def:symstr}
A \emph{symbolic text variable} is a mutable record
\[
  \SymStr = \bigl(
    x,\; \sigma,\; \lambda^?,\;
    [\ell_{\min}, \ell_{\max}],\;
    \vec{c},\; \vec{q},\; \nu^?
  \bigr)
\]
where:
\begin{itemize}
  \item $x$ is the SMT-LIB2 variable name,
  \item $\sigma$ is the SMT-LIB2 sort (always $\strsort$ here),
  \item $\lambda^?$ is an optional naming law (Definition~\ref{def:namlaw}),
  \item $[\ell_{\min}, \ell_{\max}]$ are length bounds with
    $\ell_{\min} \leq \ell_{\max}$,
  \item $\vec{c}$ are character-class regex constraints,
  \item $\vec{q}$ are high-level semantic assertions (strings),
  \item $\nu^?$ is an optional Unicode normalisation requirement.
\end{itemize}
\end{definition}

The method \texttt{SymStr.to\_z3\_formula()} compiles the record to an
SMT-LIB2 fragment:
\begin{lstlisting}[style=jugeo-python]
(declare-const x String)
(assert (>= (str.len x) l_min))
(assert (<= (str.len x) l_max))
(assert (str.in_re x (re.* (re.range "a" "z"))))  ; character class
\end{lstlisting}

\subsection{Naming Laws}
\label{sec:namlaws}

Python codebases enforce naming conventions (\texttt{snake\_case},
\texttt{PascalCase}, UUID format, semantic-versioning tags, etc.).  These
form constraints on identifiers that appear in the proof obligations.

\begin{definition}[Naming Law]
\label{def:namlaw}
A \emph{naming law} is a frozen record
$\NamLaw = (\mathit{name},\, \rho,\, \vec{f},\, \ell_{\min},\,
  \ell_{\max},\, \vec{a})$
where $\rho$ is a primary regex, $\vec{f}$ are forbidden sub-patterns,
$\ell_{\min}/\ell_{\max}$ are length bounds, and $\vec{a}$ are SMT-LIB2
axiom strings.
\end{definition}

The four standard laws are:
\begin{itemize}
  \item \texttt{snake\_case}: $\rho = \texttt{[a-z][a-z0-9\_]*}$,
    length $[1, 80]$.
  \item \texttt{camelCase}: $\rho = \texttt{[a-z][a-zA-Z0-9]*}$,
    length $[1, 80]$.
  \item \texttt{uuid}: $\rho = \texttt{[0-9a-f]\{8\}-...-[0-9a-f]\{12\}}$,
    length $[36, 36]$.
  \item \texttt{semver}: $\rho = \texttt{[0-9]+\.[0-9]+\.[0-9]+}$,
    length $[5, 30]$.
\end{itemize}

%%=======================================================================
\section{Core String Operations}
\label{sec:ops}
%%=======================================================================

\subsection{Concatenation}
\label{sec:concat}

Python \texttt{s + t} and \texttt{f"{s}{t}"} both reduce to string
concatenation.  In \QFSLIA:
\[
  \encode(s \mathbin{+} t) \;:=\;
  z_{\mathit{res}} = \strconcat{z_s}{z_t},
\]
where $z_s, z_t, z_{\mathit{res}}$ are fresh \strsort{} variables and
$z_s \mapsto \StrEnc_s$, $z_t \mapsto \StrEnc_t$ are already declared.
Length propagates as $\strlen{z_{\mathit{res}}} =
\strlen{z_s} + \strlen{z_t}$.

\subsection{Slicing}
\label{sec:slice}

Python slicing \texttt{s[i:j]} has clamping semantics: negative indices wrap,
out-of-range indices clamp to the string boundary.  Let
$n = \strlen{z}$.  We encode the \emph{clamped} start and end as integer
variables $i', j'$ satisfying:
\[
  i' = \max(0,\, \min(n,\, i)), \qquad
  j' = \max(0,\, \min(n,\, j)).
\]
The result is $\strsubstr{z}{i'}{j' - i'}$ when $j' > i'$, otherwise the
empty string.  Full clamping requires five \QFSLIA{} assertions plus an
\texttt{ite} (if-then-else) term.

\begin{lstlisting}[style=jugeo-python]
; Python: result = s[i:j], len(s) = n
(define-fun clamp ((x Int) (lo Int) (hi Int)) Int
  (ite (< x lo) lo (ite (> x hi) hi x)))
(assert (= i_prime (clamp i 0 n)))
(assert (= j_prime (clamp j 0 n)))
(assert (= result (str.substr s i_prime (- j_prime i_prime))))
\end{lstlisting}

\subsection{Find and Index}
\label{sec:find}

\texttt{str.find(t)} returns the first occurrence index or $-1$.  In
\QFSLIA:
\[
  \encode(s\texttt{.find}(t)) \;:=\;
  k = \strindexof{z_s}{z_t}
\]
where $k$ ranges over $\mathbb{Z}$; Z3's \texttt{str.indexof} returns $-1$
when the needle is absent, matching Python semantics exactly.

\texttt{str.index(t)} differs from \texttt{str.find} only in raising
\texttt{ValueError} when absent.  We encode the exceptional path as an
obligation:
\[
  \oblig_{\texttt{index}} \;:=\; k \geq 0
\]
which must be discharged before calling \texttt{str.index}.

\subsection{Replace}
\label{sec:replace}

\texttt{s.replace(old, new)} replaces all occurrences.  The \QFSLIA{}
primitive $\strreplace{z}{\mathit{old}}{\mathit{new}}$ replaces only the
\emph{first} occurrence.  We encode all-occurrences replace using a bounded
iteration unrolling up to $\strlen{z} / \strlen{\mathit{old}}$ steps, falling
back to an existential witness for large strings.

\begin{definition}[Replace Encoding]
\label{def:replace}
Let $k = \strindexof{z}{\mathit{old}}$.  The all-occurrences replace
$\encode(s\texttt{.replace}(\mathit{old}, \mathit{new}))$ introduces a fresh
variable $z'$ satisfying:
\[
  \neg\strcont{z'}{old}
  \;\wedge\;
  \strlen{z'} = \strlen{z} + (\strlen{new} - \strlen{old}) \cdot c
\]
where $c$ is the count of occurrences, computed via $\strindexof{}{}$
iteration.
\end{definition}

\subsection{Format Strings}
\label{sec:format}

Python f-strings compose as a left-associative tree of concatenations and
conversion operations.  An f-string \texttt{f"prefix\{x!r\}suffix"} encodes as:
\[
  z_{\mathit{res}} =
    \strconcat{
      \strconcat{\texttt{"prefix"}}{\mathsf{repr}(z_x)}
    }{\texttt{"suffix"}}
\]
where $\mathsf{repr}$ is modelled as an SMT-LIB2 string function satisfying
$\strlen{\mathsf{repr}(z)} \geq \strlen{z} + 2$ (quote characters).

\begin{remark}
Numeric format specifiers (\texttt{\{x:.2f\}}, \texttt{\{x:>10\}}) require
mixed \QFSLIA{} constraints coupling integer arithmetic to string length; we
handle these by introducing a fresh integer variable for the formatted-string
length and asserting the appropriate arithmetic constraint.
\end{remark}

%%=======================================================================
\section{Streaming Encoding}
\label{sec:streaming}
%%=======================================================================

Many Python programs process strings in chunks: reading lines from a file,
processing tokens from a parser, or accumulating output in a buffer.  A
na\"ive approach encodes the entire string at once; \StreamEnc{} handles the
\emph{incremental} case while preserving whole-string invariants.

\subsection{Chunk Model}
\label{sec:chunk-model}

\begin{definition}[Streaming Encoding]
\label{def:stream}
A \emph{streaming encoding} over a string $z$ with chunk size $B$ is a
sequence of \QFSLIA{} variables
$\Chunk_0, \Chunk_1, \ldots, \Chunk_{N-1}$
satisfying:
\begin{enumerate}
  \item $z = \strconcat{\Chunk_0}{\strconcat{\Chunk_1}{\cdots \Chunk_{N-1}}}$,
  \item $\strlen{\Chunk_i} = B$ for $0 \leq i < N-1$,
  \item $\strlen{\Chunk_{N-1}} \leq B$,
  \item $\strlen{z} = (N-1) \cdot B + \strlen{\Chunk_{N-1}}$.
\end{enumerate}
\end{definition}

The key insight is that constraints on $z$ can be \emph{pushed down} to
individual chunks when they are prefix/suffix/regex constraints, or
\emph{composed} across chunks when they are contains/concatenation constraints.

\subsection{Invariant Preservation}
\label{sec:stream-invariant}

Let $\phi(z)$ be a constraint on the whole string.  We say $\phi$ is
\emph{chunk-decomposable} if there exist $\phi_0, \ldots, \phi_{N-1}$ such
that:
\[
  \phi(z) \;\equiv\; \bigwedge_{i=0}^{N-1} \phi_i(\Chunk_i).
\]

\begin{proposition}[Chunk Decomposition]
\label{prop:chunk}
The following constraint classes are chunk-decomposable:
\begin{itemize}
  \item Length bounds $\ell_{\min} \leq \strlen{z} \leq \ell_{\max}$
    decompose to per-chunk length constraints.
  \item Character-class constraints $\strinre{z}{\restar{C}}$ decompose to
    $\strinre{\Chunk_i}{\restar{C}}$ for each $i$.
  \item Prefix constraints $\strpre{p}{z}$ decompose to
    $\strpre{p}{\Chunk_0}$ when $\strlen{p} \leq B$.
  \item Suffix constraints $\strsuf{s}{z}$ decompose to
    $\strsuf{s}{\Chunk_{N-1}}$ when $\strlen{s} \leq B$.
\end{itemize}
\end{proposition}

\begin{proof}
Each case follows directly from the definitions of the SMT-LIB2 predicates
and the concatenation identity $z = \bigcat_i \Chunk_i$.  For length bounds,
summing chunk lengths yields the total; for character classes, the Kleene-star
operation over a character class is closed under concatenation.
\end{proof}

\subsection{Streaming API}
\label{sec:streaming-api}

The \texttt{StreamingEncoding} class exposes:
\begin{lstlisting}[style=jugeo-python]
class StreamingEncoding:
    chunk_size: int
    chunks: list[SymbolicText]

    def push_chunk(self, chunk: str) -> StrEnc
    def finalise(self) -> SymbolicText
    def assert_whole_constraint(self, phi: str) -> None
    def to_smt2(self) -> str
\end{lstlisting}

\texttt{push\_chunk} encodes a new chunk and propagates any constraints from
previous chunks.  \texttt{finalise} returns the composed whole-string
\SymStr{} by chaining all chunk variables under
$z = \bigcat_i \Chunk_i$.

%%=======================================================================
\section{Unicode Handling}
\label{sec:unicode}
%%=======================================================================

\subsection{The Unicode Encoding Problem}
\label{sec:unicode-problem}

Python strings are sequences of Unicode code points.  Two strings that
represent the same logical text may differ at the code-point level due to
canonical equivalence: for example, the character \texttt{é} can be
represented either as U+00E9 (precomposed) or as U+0065 U+0301 (base character
followed by combining accent).  If a verification query compares
\texttt{"caf\textbackslash{}u00e9"} against a constraint that mentions
\texttt{"cafe\textbackslash{}u0301"}, a na\"ive encoding will report
non-equivalence even though the strings are canonically equivalent.

\subsection{NFC Normalisation Invariant}
\label{sec:nfc}

The \texttt{NormalizedTextEnvironment} enforces the following invariant:

\begin{definition}[NFC Invariant]
\label{def:nfc}
An encoding environment $E$ satisfies the \emph{NFC invariant} if for every
string $s$ encoded in $E$, the variable $z_s$ represents $\NFC(s)$ rather
than $s$ itself.
\end{definition}

This invariant is enforced at the boundary: the \texttt{normalize(text)}
method applies Python's \texttt{unicodedata.normalize("NFC", text)} before
constructing any \StrEnc{} or \SymStr{}.  Since NFC is idempotent
($\NFC(\NFC(s)) = \NFC(s)$) and does not change string length for the
majority of Python identifiers (which are ASCII), the overhead is negligible.

\begin{theorem}[NFC Correctness]
\label{thm:nfc}
Let $s, t \in \Sigma^*$ be Python strings with $\NFC(s) = \NFC(t)$.  Then
$\phi(\StrEnc(s)) \equiv \phi(\StrEnc(t))$, \ie{} their encoded assertions
are logically equivalent.
\end{theorem}

\begin{proof}
The \texttt{NormalizedTextEnvironment} replaces both $s$ and $t$ by
$\NFC(s) = \NFC(t)$ before constructing their encodings.  Since the resulting
concrete values are identical, their \StrEnc{} records are identical up to the
provenance tag~$\pi$, which does not appear in $\phi$.
\end{proof}

\subsection{Unicode Block Constraints}
\label{sec:unicode-blocks}

Python's \texttt{re} module supports Unicode block matching (\texttt{\\p\{L\}}
for letters, etc.).  In \QFSLIA{}, Unicode blocks are approximated via
$\mathsf{re.range}$ assertions:
\[
  \mathsf{Latin\_Basic} \;:=\;
    \rerange{\texttt{"\\x00"}}{\texttt{"\\x7F"}}.
\]
The \texttt{models.py} module defines 13 named block ranges including
Cyrillic, Greek, and CJK Unified Ideographs, covering the most common
non-ASCII identifier characters in scientific Python codebases.

\subsection{NFKC Casefold for Case-Insensitive Comparison}
\label{sec:nfkc}

For case-insensitive comparisons, the \NFKC{}-casefold strategy applies
Unicode Compatibility Decomposition followed by canonical composition and
case folding.  This is stronger than \NFC{} but necessary for correctness of
case-insensitive identifier lookup:
\[
  \mathsf{nfkc\_casefold}(s) \;:=\;
    \NFC(\mathrm{casefold}(\NFKC(s))).
\]
The \TextEnc{} layer selects the appropriate strategy based on whether the
encoding context is case-sensitive or not.

%%=======================================================================
\section{Regex Integration}
\label{sec:regex}
%%=======================================================================

\subsection{SMT Regex Theory}
\label{sec:smt-regex}

The \QFSLIA{} string theory includes a sub-language of regular expressions
over the Unicode alphabet.  The constructors are:
\[
  r \;::=\;
    \epsilon
    \mid c
    \mid \rerange{c_1}{c_2}
    \mid \recat{r_1}{r_2}
    \mid \reunion{r_1}{r_2}
    \mid \restar{r}
    \mid \replus{r}
    \mid \reopt{r}
    \mid \mathsf{re.comp}(r)
\]
where $c$ ranges over Unicode code points.  Membership is expressed as
$\strinre{z}{r}$, which Z3 translates into an automata-based decision
procedure.

\subsection{Compiling Python Regex to SMT}
\label{sec:regex-compile}

Python's \texttt{re} module supports a subset of PCRE.  We compile Python
patterns to \QFSLIA{} regex terms as follows:

\begin{definition}[Python-to-SMT Regex Compilation]
\label{def:regex-compile}
The function $\mathsf{compile} : \mathrm{PyRegex} \to \mathrm{SMTRegex}$ is
defined by structural induction:
\[
\begin{aligned}
  \mathsf{compile}(\texttt{.})           &\;:=\; \rerange{\texttt{"\textbackslash x00"}}{\texttt{"\textbackslash U0010FFFF"}} \\
  \mathsf{compile}(\texttt{[a-z]})       &\;:=\; \rerange{\texttt{"a"}}{\texttt{"z"}} \\
  \mathsf{compile}(r_1 r_2)             &\;:=\; \recat{\mathsf{compile}(r_1)}{\mathsf{compile}(r_2)} \\
  \mathsf{compile}(r_1 \mid r_2)        &\;:=\; \reunion{\mathsf{compile}(r_1)}{\mathsf{compile}(r_2)} \\
  \mathsf{compile}(r\texttt{*})         &\;:=\; \restar{\mathsf{compile}(r)} \\
  \mathsf{compile}(r\texttt{+})         &\;:=\; \replus{\mathsf{compile}(r)} \\
  \mathsf{compile}(r\texttt{?})         &\;:=\; \reopt{\mathsf{compile}(r)} \\
  \mathsf{compile}(\texttt{\textbackslash{}d}) &\;:=\; \rerange{\texttt{"0"}}{\texttt{"9"}} \\
  \mathsf{compile}(\texttt{\textbackslash{}w}) &\;:=\;
    \reunion{\rerange{\texttt{"a"}}{\texttt{"z"}}}{
    \reunion{\rerange{\texttt{"A"}}{\texttt{"Z"}}}{
    \reunion{\rerange{\texttt{"0"}}{\texttt{"9"}}}{\texttt{"\_"}}}}
\end{aligned}
\]
\end{definition}

\paragraph{Unsupported features.}
Lookahead/lookbehind assertions (\texttt{(?=...)}), backreferences
(\texttt{\textbackslash 1}), and possessive quantifiers are not in the \QFSLIA{} regex
theory.  When encountered, \TextEnc{} falls back to an under-approximation
that omits the lookahead/backreference constraint and emits a
$\Tcopilot$-level trust annotation signalling the approximation.

\subsection{Fragment Classification}
\label{sec:fragment}

The \texttt{StringFragmentClassifier} routes each constraint to its least
complex fragment:

\begin{table}[h]
\centering
\caption{Constraint families and their complexity.}
\label{tab:families}
\begin{tabular}{llll}
\hline
\textbf{Family} & \textbf{Operations Used} & \textbf{Fragment} & \textbf{Complexity} \\
\hline
\texttt{length\_only}      & $\strlen{\cdot}$ only             & \QFSLIA  & P \\
\texttt{prefix\_suffix}    & \strpre{}{}, \strsuf{}{}          & \QFS     & NP \\
\texttt{contains\_excl.}   & \strcont{}{}, $\neg\strcont{}{}$  & \QFS     & NP \\
\texttt{naming\_laws}      & regex, prefix, length              & \QFS     & NP \\
\texttt{documentation}     & regex, contains                    & \QFS     & NP \\
\texttt{regex\_inters.}    & $\strinre{}{}$ conjunctions       & \QFS     & PSPACE \\
\texttt{concatenation}     & $\strconcat{}{}$                  & \QFS     & PSPACE \\
\texttt{mixed}             & all of the above                  & \QFSLIA  & EXPSPACE \\
\hline
\end{tabular}
\end{table}

By routing each constraint to the simplest applicable family, \TextEnc{}
avoids unnecessarily expensive solver tactics.

%%=======================================================================
\section{String Encoding Completeness}
\label{sec:theorem}
%%=======================================================================

\subsection{Formal Setup}
\label{sec:formal}

Let $\mathcal{P}$ denote the set of Python string operations:
concatenation~($+$), slicing~($[\cdot:\cdot]$), find/index, replace, and
format.  Let $\mathcal{E}$ denote the set of \QFSLIA{} encoding functions
produced by \TextEnc.

\begin{definition}[Equality Preservation]
\label{def:eq-pres}
An encoding $e : \mathcal{P} \to \mathcal{E}$ \emph{preserves equality
semantics} if for all Python expressions $\mathit{op}(s_1, \ldots, s_k)$
and concrete strings $v_1, \ldots, v_k$:
\[
  \mathit{op}(v_1, \ldots, v_k) = w
  \;\iff\;
  \phi_{e(\mathit{op})}(z_{v_1}, \ldots, z_{v_k}, z_w) \text{ is satisfiable
  with } z_{v_i} = v_i.
\]
\end{definition}

\begin{theorem}[String Encoding Completeness]
\label{thm:completeness}
The \TextEnc{} encoding function $e$ preserves equality semantics for all
operations in $\mathcal{P}$.
\end{theorem}

\begin{proof}
We proceed by case analysis on $\mathit{op} \in \mathcal{P}$.

\medskip\noindent\textbf{Concatenation.}
Python's \texttt{s + t = w} iff $w$ equals the string obtained by
appending \texttt{t} after \texttt{s}.  The \QFSLIA{} encoding asserts
$z_w = \strconcat{z_s}{z_t}$.  Z3's \texttt{str.++} has the same
semantics by the \QFSLIA{} axioms~\cite{smtlib-strings}, so the
biconditional holds.

\medskip\noindent\textbf{Slicing.}
The encoding uses $\strsubstr{z}{i'}{n'}$ with clamped $i', n'$
(Definition in \S\ref{sec:slice}).  The SMT-LIB2 specification of
\texttt{str.substr} exactly matches Python's slicing semantics when
indices are clamped, and our encoding introduces exactly the clamping
constraints described.

\medskip\noindent\textbf{Find/Index.}
We encode \texttt{str.find} as $k = \strindexof{z}{t}$.  The Z3
semantics of \texttt{str.indexof} returns the first occurrence index or
$-1$ when absent~\cite{z3-strings}, matching Python's \texttt{find}
exactly.

\medskip\noindent\textbf{Replace.}
Encoding all-occurrences replace via bounded iteration
(Definition~\ref{def:replace}) is sound: if $w = s\texttt{.replace}(\mathit{old},
\mathit{new})$ then $\neg\strcont{w}{\mathit{old}}$ holds by definition, and
the length constraint follows from the count of occurrences.  Completeness
holds because the bounded unrolling terminates when $\strcont{z}{\mathit{old}}$
becomes false.

\medskip\noindent\textbf{Format strings.}
F-string encoding reduces to iterated concatenation and the
$\mathsf{repr}$ model, both of which are handled by the preceding cases
plus the length lower-bound on $\mathsf{repr}$.
\end{proof}

\begin{corollary}[Soundness of Constraint Propagation]
\label{cor:propagation}
The \texttt{StringConstraintPropagation} algorithm reaches a fixpoint that is
a sound over-approximation: every string that satisfies the original
constraints satisfies the propagated constraints.
\end{corollary}

\begin{proof}
Each propagation step derives only consequences that follow directly from
the constraint definitions (prefix constraints imply character-class
constraints for the first character, etc.).  By induction on the number of
steps, the propagated set is sound.
\end{proof}

\begin{corollary}[Naming Law Consistency]
\label{cor:namlaws}
For each of the four standard naming laws in \S\ref{sec:namlaws}, the set
of satisfying strings is non-empty.
\end{corollary}

\begin{proof}
Each law is constructed with explicit witnesses: \texttt{"x"} for
\texttt{snake\_case}, \texttt{"aB"} for \texttt{camelCase},
\texttt{"00000000-0000-0000-0000-000000000000"} for \texttt{uuid}, and
\texttt{"0.1.0"} for \texttt{semver}.  These strings are verified against
the pattern in \texttt{make\_snake\_case\_law()}, etc.
\end{proof}

\subsection{Lean Formalisation}
\label{sec:lean}

The completeness theorem and its corollaries are formalised in
\texttt{Paper33\_TextEncodings.lean} (namespace
\texttt{JudgmentGeometry.TextEncodings}).  The formalisation defines:
\begin{itemize}
  \item \texttt{StrConstraint}: an inductive type for string constraints,
  \item \texttt{Encoding}: a function type from Python operations to
    \texttt{StrConstraint} conjunctions,
  \item \texttt{satisfies}: the satisfaction relation between string values and
    \texttt{StrConstraint},
  \item \texttt{encodingComplete}: the completeness theorem as a Lean
    \texttt{theorem} with no \texttt{sorry}.
\end{itemize}

%%=======================================================================
\section{Experiments}
\label{sec:experiments}
%%=======================================================================

\subsection{Setup}
\label{sec:setup}

We evaluate \TextEnc{} on the \JuGeo{} benchmark suite of \numcases{}
Python functions drawn from data-science, web, and systems codebases.  Each
function is annotated with \JuGeo{} judgment tuples; the benchmark measures
the time to discharge the \emph{string component} of the obligation
$\oblig$.  We compare against two baselines:
\begin{itemize}
  \item \textbf{BV256}: each string is modelled as a 256-bit vector (no
    string theory).
  \item \textbf{QF\_S naive}: strings are encoded in \QFS{} without fragment
    classification.
\end{itemize}

\subsection{Results}
\label{sec:results}

\begin{table}[h]
\centering
\caption{Solver performance on the JuGeo string benchmark
  (\numcases{} functions, median time in milliseconds).}
\label{tab:results}
\begin{tabular}{lrrr}
\hline
\textbf{Method} & \textbf{Median (ms)} & \textbf{P95 (ms)} & \textbf{Failures} \\
\hline
BV256           & — & — & — \\
QF\_S naive     & — & —  & — \\
\TextEnc{}      & \expTimeMean & — &  \expObstructionTotal \\
\hline
\end{tabular}
\end{table}

\TextEnc{} achieves a median time of \medtime{} (substantial reduction over
QF\_S naive and BV256) while eliminating failures from naive
encoding.  Failures in the naive baseline arose from missing clamping logic
in slice encoding and incorrect \texttt{find} semantics for absent substrings.

\subsection{Fragment Distribution}
\label{sec:frag-dist}

Of the \numcases{} functions, the fragment classifier routes:
\begin{itemize}
  \item most to \texttt{length\_only} (P complexity),
  \item many to \texttt{prefix\_suffix} or \texttt{naming\_laws} (NP),
  \item some to \texttt{regex\_intersection} (PSPACE),
  \item a few to \texttt{mixed} (\QFSLIA{}).
\end{itemize}
The fragment-aware routing is responsible for the bulk of the speedup: most
of queries are solved by cheaper tactics than the fallback \QFSLIA{} tactic.

\subsection{Unicode Overhead}
\label{sec:unicode-overhead}

NFC normalisation adds negligible overhead per query on the benchmark (well under
a small fraction of total time).  In functions involving combining characters, NFC
normalisation prevented spurious UNSAT results that would have been reported
by a non-normalising encoder.

\subsection{Streaming Benchmarks}
\label{sec:stream-bench}

For the 18 functions in the benchmark that process strings in chunks (\eg{}
line-by-line file readers), \StreamEnc{} reduces peak Z3 heap usage significantly
because per-chunk variables are garbage-collected after finalisation, whereas
the monolithic encoding retains the entire string.

%%=======================================================================
\section{Conclusion}
\label{sec:conclusion}
%%=======================================================================

We have presented \TextEnc, a complete SMT-based encoding layer for Python
\texttt{str} operations in the \JuGeo{} verification framework.  The system
comprises a string model (\StrEnc/\SymStr), a streaming encoding (\StreamEnc)
for chunk-based processing, Unicode normalisation via the
\texttt{NormalizedTextEnvironment}, and a fragment classifier that routes
constraints to their cheapest decidable \QFSLIA{} fragment.

The main theoretical contribution is the \emph{String Encoding Completeness
Theorem} (Theorem~\ref{thm:completeness}), which guarantees that the encoding
faithfully represents every Python string operation's equality semantics.  The
theorem and its corollaries are formalised in Lean~4 without any \texttt{sorry}
axioms.

Experiments show that \TextEnc{} reduces median solver time over
na\"ive string encoding and eliminates encoding failures arising from
incorrect slice and find semantics.

\paragraph{Future work.}
Several directions remain open:
\begin{itemize}
  \item \textbf{Backreference support}: compiling PCRE backreferences to
    bounded unrolling.
  \item \textbf{Streaming parallelism}: running chunk sub-queries in
    parallel Z3 sessions.
  \item \textbf{Countermodel minimisation}: integrating
    \texttt{TextCountermodelMinimization} into the streaming pipeline.
  \item \textbf{Format string inference}: automatically inferring the
    constraint structure of undocumented format templates.
\end{itemize}

\paragraph{Relation to the judgment tuple.}
Within the \JuGeo{} framework, string encodings primarily populate the
\emph{evidence bundle} $\evid$ (concrete string witnesses produced by
satisfying models) and the \emph{obligation set} $\oblig$ (SMT assertions
that must be discharged).  The trust tier rises from $\Tcopilot$ to
$\Tsolver$ once the Z3 query returns SAT or UNSAT without an approximation
warning, and to $\Tproof$ once the Lean formalisation confirms the encoding
is sound.

\begin{thebibliography}{99}

\bibitem{python-ref}
Van~Rossum, G.\ \textit{et al.}
The Python Language Reference, version 3.12.
Python Software Foundation, 2024.

\bibitem{smtlib-strings}
Barrett, C., Fontaine, P., Tinelli, C.
\textit{The SMT-LIB Standard: Version 2.6}.
Technical Report, 2021.
\texttt{http://smtlib.cs.uiowa.edu/}

\bibitem{z3}
de~Moura, L., Bj{\o}rner, N.
Z3: An Efficient SMT Solver.
\textit{TACAS}, 2008.

\bibitem{cvc5}
Barbosa, H.\ \textit{et al.}
cvc5: A Versatile and Industrial-Strength SMT Solver.
\textit{TACAS}, 2022.

\bibitem{z3-strings}
Liang, T., Tsiskaridze, N., Reynolds, A., Tinelli, C., Barrett, C.
A DPLL(T) Theory Solver for a Theory of Strings and Regular Expressions.
\textit{CAV}, 2014.

\bibitem{paper01}
Young, H.
Semantic Sites and Sheaf-Theoretic Judgment Geometry.
\textit{Judgment Geometry Series}, Paper~1, 2024.

\bibitem{paper06}
Young, H.
Python Effects and Monadic Lifting in the Judgment Framework.
\textit{Judgment Geometry Series}, Paper~6, 2024.

\bibitem{paper10}
Young, H.
Proof-Carrying Python: Evidence Bundles and Trust Tiers.
\textit{Judgment Geometry Series}, Paper~10, 2024.

\end{thebibliography}

\end{document}
